\section{Risultati}

\vspace*{\fill}
\subsection{Prestazioni Generali dei Modelli}
\subsubsection{Tasso di Successo e Tecniche di Quantizzazione}

La valutazione delle prestazioni dei modelli è stata condotta su una workstation con le seguenti specifiche tecniche:

\begin{itemize}
    \item CPU: \textbf{AMD Ryzen 7 5800x}
    \item RAM: \textbf{32GB}
    \item GPU: \textbf{AMD Radeon RX 7800XT con 16GB di VRAM}
    \item Sistema Operativo: \textbf{Windows 11 23H2}
\end{itemize}
Nell'ambito di questo studio, cinque dei sei modelli esaminati hanno dimostrato dimensioni sufficientemente contenute da permettere un'esecuzione completa sulla GPU senza incontrare difficoltà significative. Tuttavia, il modello \textbf{Gemma-2-27b-Q3\_K\_S} ha rappresentato un'eccezione: a causa delle sue dimensioni considerevoli, non è stato possibile delegare l'intera inferenza alla GPU, comportando un notevole incremento dei tempi di inferenza.
È opportuno sottolineare che, al momento di questa ricerca, il supporto di \textbf{ROCm} su \textbf{Windows} per le GPU \textbf{AMD} è ancora "acerbo" e non ha raggiunto una piena stabilità.
Pertanto, le prestazioni osservate non riflettono lo stato dell'arte dell'esperienza che si può ottenere con LLM self-hosted su hardware \textbf{NVIDIA} nella fascia consumer.
Questa limitazione tecnica deve essere considerata nell'interpretazione dei risultati ottenuti e nel confronto con altre piattaforme hardware.
\vspace*{\fill}

\clearpage


\vspace*{\fill}
\subsection{Tasso di Successo}
Il tasso di successo, definito come la percentuale di interazioni completate senza errori significativi o incongruenze, varia tra i modelli testati
\begin{figure}[H]
    \centering
    \centering
    \includegraphics[width=\textwidth]{success_rate_plot}
    \caption{Tasso di successo per modello}
    \label{fig:success_rate_plot}
\end{figure}
\paragraph{Analisi}
I tassi di successo variano da \textbf{95\%} per il modello meno performante a \textbf{100\%} per il più performante. In particolare, i modelli \textbf{Gemma} mostrano tassi di successo superiori (\textbf{99\%}-\textbf{100\%}) rispetto ai modelli \textbf{LLaMA} (\textbf{95\%}-\textbf{97\%}).
Questo suggerisce una maggiore affidabilità e coerenza nei risultati ottenuti dai modelli \textbf{Gemma}, specialmente quelli da \textbf{27B}, rispetto ai modelli LLaMA.\\

\vspace*{\fill}
\clearpage


\vspace*{\fill}
\subsection{Accuratezza per singolo tratto}
\begin{figure}[H]
    \centering
    \includegraphics[width=\textwidth]{trait_accuracy_boxplot}
    \caption{Distribuzione dell'accuratezza per tratto di personalità e modello}
    \label{fig:trait_accuracy_boxplot}
\end{figure}
Si nota una consistente superiorità dei modelli \textbf{Gemma}, in particolare per il \textbf{Nevroticismo} e la \textbf{Coscienziosità}. L'\textbf{Estroversione}, invece, emerge come il tratto più difficile da predire accuratamente per tutti i modelli.
\paragraph{Analisi}
L'Apertura emerge tra i più difficili da predire, con accuratezze medie che variano da \textbf{55\%} a \textbf{95\%} tra i modelli.
Questa difficoltà potrebbe essere dovuta alla natura complessa e sfaccettata di questo tratto, che richiede una maggiore sensibilità e comprensione delle sfumature del linguaggio.\\
D'altra parte, il \textbf{Nevroticismo} e la \textbf{Coscienziosità} sono meglio predetti dai modelli, con accuratezze medie rispettivamente del \textbf{92\%} e \textbf{82\%}. Ciò potrebbe indicare che questi tratti hanno manifestazioni più coerenti e riconoscibili nel linguaggio.


\vspace*{\fill}
\clearpage

\vspace*{\fill}
\subsection{Scatter Plot Durata vs Accuratezza}
\begin{figure}[H]
    \centering
    \includegraphics[width=\textwidth]{duration_accuracy_scatter}
    \caption{Durata vs Accuratezza Totale per Modello}
    \label{fig:duration_accuracy_scatter}
\end{figure}

\paragraph{Trade-off tra accuratezza e tempo di inferenza}

Il grafico evidenzia un chiaro compromesso tra accuratezza e tempo di inferenza per diversi modelli di linguaggio.

\begin{table}[h]
\centering
\begin{tabular}{|l|c|c|}
\hline
\textbf{Modello} & \textbf{Accuratezza (\%)} & \textbf{Tempo di inferenza (s)} \\
\hline
LLaMA & 55-92 & 5-7 \\
Gemma 9B & 68-92 & 10-15 \\
Gemma 27B & 72-92 & 15-20 (IQ2\_M) / 30-80 (Q3\_K\_S) \\
\hline
\end{tabular}
\end{table}
I modelli \textbf{Gemma} offrono maggiore accuratezza e coerenza, ma con tempi di inferenza più lunghi, mentre i modelli \textbf{LLaMA} privilegiano la velocità a scapito di una leggera perdita di coerenza nell'accuratezza.

\begin{itemize}
    \item \textbf{LLaMA 3.1 8B}: Massima velocità, Scarsa accuratezza
    \item \textbf{Gemma 2 9B}: Velocità accettabile, Maggiore accuratezza
    \item \textbf{Gemma 2 27B}: Estremamente lento, Massima accuratezza
\end{itemize}

La scelta del modello dipende dall'applicazione:
\begin{itemize}
    \item Per utilizzo in tempo reale o su grandi contesti: preferibili i modelli \textbf{LLaMA}
    \item Per massima precisione: più adatti i modelli \textbf{Gemma}
\end{itemize}

\vspace*{\fill}
\clearpage

\vspace*{\fill}
\subsection{Consistenza delle accuratezze per modello}
\begin{figure}[H]
    \centering
    \includegraphics[width=\textwidth]{accuracy_boxplot}
    \caption{Distribuzione delle accuratezze totali per modello}
    \label{fig:accuracy_boxplot}
\end{figure}
I modelli \textbf{Gemma} mostrano mediane e medie di accuratezza più elevate rispetto ai modelli
\textbf{LLaMA}, con \textbf{Gemma-2-27b-IQ2\_M} che raggiunge le prestazioni migliori. Le tecniche di quantizzazione
sembrano avere un impatto minimo sulle prestazioni complessive dei modelli \textbf{Gemma}.\\\\

\vspace*{\fill}
\clearpage

\vspace*{\fill}
\subsection{Istogrammi delle Accuratezze per Tratto}

L'analisi delle distribuzioni di accuratezza per ciascun tratto di personalità fornisce ulteriori approfondimenti sulle prestazioni dei modelli testati.

\subsubsection{Gradevolezza (Agreeableness)}

\begin{figure}[h]
    \centering
    \includegraphics[width=\textwidth]{histogram_Agreeableness}
    \caption{Distribuzione dell'accuratezza per Gradevolezza}
    \label{fig:accuracy_agreeableness}
\end{figure}
L'istogramma relativo alla \textbf{Gradevolezza} mostra una distribuzione generalmente più ampia rispetto agli altri tratti. I modelli \textbf{Gemma} si distinguono per una lieve superiorità, con distribuzioni che tendono a concentrarsi su valori di accuratezza più elevati. Questa maggiore ampiezza suggerisce che la Gradevolezza sia un tratto con manifestazioni linguistiche più variegate, rendendo la sua previsione più complessa, ma comunque affrontabile dai modelli più avanzati.

\vspace*{\fill}
\clearpage

\vspace*{\fill}
\subsubsection{Coscienziosità (Conscientiousness)}

\begin{figure}[h]
    \centering
    \includegraphics[width=\textwidth]{histogram_Conscientiousness}
    \caption{Distribuzione dell'accuratezza per Coscienziosità}
    \label{fig:accuracy_conscientiousness}
\end{figure}
Per la \textbf{Coscienziosità}, i modelli \textbf{Gemma} mostrano distribuzioni più concentrate verso valori di accuratezza più alti rispetto ai modelli \textbf{LLaMA}.
Questo suggerisce una migliore capacità dei modelli Gemma nel cogliere i segnali linguistici associati alla \textbf{Coscienziosità}, forse grazie a una più sofisticata comprensione delle sfumature del linguaggio relative all'organizzazione, alla disciplina e all'attenzione ai dettagli.

\vspace*{\fill}
\clearpage

\vspace*{\fill}
\subsubsection{Estroversione (Extraversion)}

\begin{figure}[h]
    \centering
    \includegraphics[width=\textwidth]{histogram_Extraversion}
    \caption{Distribuzione dell'accuratezza per Estroversione}
    \label{fig:accuracy_extraversion}
\end{figure}
L'istogramma dell'\textbf{Estroversione} rivela le distribuzioni più variabili e disperse tra tutti i tratti.
Questa caratteristica conferma la difficoltà nella previsione accurata dell'\textbf{Estroversione} per tutti i modelli. La variabilità potrebbe essere attribuita alla complessità di questo tratto, che può manifestarsi in modi diversi nel linguaggio scritto, rendendo la sua identificazione più impegnativa per i modelli di linguaggio.

\vspace*{\fill}
\clearpage

\vspace*{\fill}
\subsubsection{Nevroticismo (Neuroticism)}

\begin{figure}[h]
    \centering
    \includegraphics[width=\textwidth]{histogram_Neuroticism}
    \caption{Distribuzione dell'accuratezza per Nevroticismo}
    \label{fig:accuracy_neuroticism}
\end{figure}
Per il \textbf{Nevroticismo}, tutti i modelli mostrano una distribuzione dell'accuratezza spostata verso destra, indicando prestazioni generalmente buone. I modelli \textbf{Gemma}, in particolare quelli da \textbf{27B}, mostrano distribuzioni più concentrate verso valori di accuratezza più alti. Questo suggerisce che il \textbf{Nevroticismo} potrebbe avere manifestazioni linguistiche più facilmente riconoscibili, permettendo ai modelli di identificarlo con maggiore precisione.

\vspace*{\fill}
\clearpage

\vspace*{\fill}
\subsubsection{Apertura (Openness)}

\begin{figure}[h]
    \centering
    \includegraphics[width=\textwidth]{histogram_Openness}
    \caption{Distribuzione dell'accuratezza per Apertura}
    \label{fig:accuracy_openness}
\end{figure}

Le distribuzioni per l'Apertura (Figura \ref{fig:accuracy_openness}) sono più variabili tra i modelli. I modelli Gemma da 27B mostrano distribuzioni più ampie, suggerendo una maggiore capacità di catturare le sfumature di questo tratto. L'Apertura, essendo un tratto legato alla creatività e alla curiosità intellettuale, potrebbe manifestarsi in modi diversi nel linguaggio, richiedendo modelli più sofisticati per una previsione accurata.

In generale, i modelli Gemma, specialmente quelli da 27B, mostrano distribuzioni di accuratezza più favorevoli per la maggior parte dei tratti. Questo conferma le osservazioni precedenti sulla loro superiore capacità di catturare e prevedere le sfumature della personalità. La variabilità nelle distribuzioni tra i diversi tratti sottolinea l'importanza di considerare le prestazioni specifiche per ciascun tratto nella valutazione complessiva dei modelli, piuttosto che basarsi solo su metriche aggregate.

\vspace*{\fill}
\clearpage

\subsection{Heatmap delle Correlazioni}

\subsubsection{Gemma-2-27b-IQ2\_M}

\begin{figure}[h]
    \centering
    \includegraphics[width=\textwidth]{heatmap_Gemma-2-27b-IQ2_M}
    \caption{Correlazione tra valori attesi e ottenuti - Gemma-2-27b-IQ2\_M}
    \label{fig:heatmap_Gemma-2-27b-IQ2_M}
\end{figure}

\textbf{Gemma-2-27b-IQ2\_M}: Questo modello esibisce correlazioni molto elevate per la maggior parte dei tratti. Il \textbf{Nevroticismo} emerge come il tratto meglio replicato, con una correlazione di \textbf{0.81}. La \textbf{Coscienziosità} (\textbf{0.57}) e la \textbf{Gradevolezza} (\textbf{0.55}) mostrano correlazioni moderate, mentre gli altri tratti presentano una precisione inferiore.

\vspace*{\fill}

\clearpage

\vspace*{\fill}

\subsubsection{Gemma-2-27b-Q3\_K\_S}

\begin{figure}[h]
    \centering
    \includegraphics[width=\textwidth]{heatmap_Gemma-2-27b-Q3_K_S}
    \caption{Correlazione tra valori attesi e ottenuti - Gemma-2-27b-Q3\_K\_S}
    \label{fig:heatmap_Gemma-2-27b-Q3_K_S}
\end{figure}

\textbf{Gemma-2-27b-Q3\_K\_S}: Questo modello mostra prestazioni analoghe al più quantizzato \textbf{IQ2\_M}, ma con corrispondenze notevolmente più elevate per alcuni tratti.
In particolare, il \textbf{Nevroticismo} (\textbf{0.91}), \textbf{l'Apertura} (\textbf{0.72}) e la \textbf{Coscienziosità} (\textbf{0.70}) presentano correlazioni piuttosto alte. Anche gli altri tratti mostrano correlazioni più soddisfacenti rispetto al modello precedente.

\vspace*{\fill}

\clearpage

\vspace*{\fill}

\subsubsection{Gemma-2-9b-Q5\_K\_M}

\begin{figure}[h]
    \centering
    \includegraphics[width=\textwidth]{heatmap_Gemma-2-9b-Q5_K_M}
    \caption{Correlazione tra valori attesi e ottenuti - Gemma-2-9b-Q5\_K\_M}
    \label{fig:heatmap_Gemma-2-9b-Q5_K_M}
\end{figure}

\textbf{Gemma-2-9b-Q5\_K\_M}: Questo modello presenta un pattern di correlazioni simile al modello \textbf{Q8\_0}, ma con valori leggermente inferiori per la maggior parte dei tratti. Il Nevroticismo si conferma il tratto meglio replicato, con una correlazione di \textbf{0.91}.

\vspace*{\fill}

\clearpage

\vspace*{\fill}

\subsubsection{Gemma-2-9b-Q8\_0}
\begin{figure}[h]
    \centering
    \includegraphics[width=\textwidth]{heatmap_Gemma-2-9b-Q8_0}
    \caption{Correlazione tra valori attesi e ottenuti - Gemma-2-9b-Q8\_0}
    \label{fig:heatmap_Gemma-2-9b-Q8_0}
\end{figure}

\textbf{Gemma-2-9b-Q8\_0}: Questo modello dimostra prestazioni eccellenti nella replicazione del Nevroticismo, con una correlazione di \textbf{0.91}. Presenta inoltre buone correlazioni per l'\textbf{Estroversione} (\textbf{0.66}) e la \textbf{Coscienziosità} (\textbf{0.46}). L'\textbf{Apertura} (\textbf{0.51}) e la \textbf{Gradevolezza} (\textbf{0.57}) mostrano correlazioni moderate.

\vspace*{\fill}

\clearpage

\vspace*{\fill}

\subsubsection{LLaMA-3.1-8B-Q5\_K\_M}

\begin{figure}[h]
    \centering
    \includegraphics[width=\textwidth]{heatmap_LLaMA-3.1-8B-Q5_K_M}
    \caption{Correlazione tra valori attesi e ottenuti - LLaMA-3.1-8B-Q5\_K\_M}
    \label{fig:heatmap_LLaMA-3.1-8B-Q5_K_M}
\end{figure}

\textbf{LLaMA-3.1-8B-Q5\_K\_M}: Questo modello mostra correlazioni generalmente più basse rispetto ai modelli \textbf{Gemma}. Il tratto replicato meglio è la \textbf{Coscienziosità} con una correlazione di \textbf{0.60}. Il \textbf{Nevroticismo} presenta una correlazione di \textbf{0.54}, mentre gli altri tratti mostrano correlazioni inferiori.

\vspace*{\fill}

\clearpage

\vspace*{\fill}

\subsubsection{LLaMA-3.1-8B-Q8\_0}

\begin{figure}[h]
    \centering
    \includegraphics[width=\textwidth]{heatmap_LLaMA-3.1-8B-Q8_0}
    \caption{Correlazione tra valori attesi e ottenuti - LLaMA-3.1-8B-Q8\_0}
    \label{fig:heatmap_LLaMA-3.1-8B-Q8_0}
\end{figure}

\textbf{LLaMA-3.1-8B-Q8\_0}: Migliora leggermente rispetto al modello \textbf{Q5\_K\_M}, con correlazioni più elevate per la \textbf{Coscienziosità} (\textbf{0.69}). Tuttavia, le correlazioni rimangono generalmente più basse rispetto ai modelli \textbf{Gemma}.

\vspace*{\fill}

\clearpage


\subsection{Analisi Complessiva delle Prestazioni dei Modelli}

Emerge chiaramente la superiorità dei modelli \textbf{Gemma} nella replicazione accurata di tutti i tratti, con prestazioni particolarmente notevoli per il \textbf{Nevroticismo} e la \textbf{Coscienziosità}.\\\\
I \textbf{modelli Gemma 27B} mostrano generalmente prestazioni superiori rispetto alle loro controparti \textbf{9B}, suggerendo che un modello più grande, anche se con una quantizzazione più aggressiva può fornire migliori capacità di comprensione del contesto.
Questo potrebbe essere attribuito alla maggiore capacità di questi modelli di cogliere sfumature sottili e complesse del linguaggio associato ai vari tratti di personalità.\\\\
I \textbf{modelli LLaMA}, sebbene generalmente accurati, mostrano correlazioni più basse rispetto ai modelli \textbf{Gemma} per la maggior parte dei tratti.
Ciò potrebbe indicare una minor capacità di catturare le sfumature più sottili dei vari tratti di personalità.
Tuttavia, è importante notare che questi modelli offrono un enorme vantaggio in termini di velocità di inferenza, che li rendono assolutamente preferibili in scenari dove la rapidità di elaborazione è cruciale.\\\\
L'\textbf{Estroversione} si conferma come il tratto più sfidante da replicare in modo consistente per tutti i modelli.
Questa difficoltà potrebbe essere dovuta alla natura complessa di questo tratto, che può manifestarsi in modi diversi nel linguaggio scritto, al contrario, il \textbf{Nevroticismo} emerge come il tratto più facilmente replicabile, con correlazioni molto alte per la maggior parte dei modelli. Questo suggerisce che il linguaggio associato al Nevroticismo potrebbe avere caratteristiche più distintive e facilmente riconoscibili per i modelli di linguaggio.

\subsection{Impatto della Quantizzazione sulle Prestazioni}

L'analisi delle heatmap rivela anche l'impatto delle diverse tecniche di quantizzazione sulle prestazioni dei modelli:

\begin{itemize}
    \item La quantizzazione a \textbf{3 bit (Q3\_K\_S)} nel modello \textbf{Gemma-2-27b} offre le performance migliori in termini di accuratezza. Tuttavia, questo viene a discapito di una maggiore dimensione del modello, che in questo caso si traduce in un overload di risorse computazionali.

    \item La quantizzazione a \textbf{2 bit (IQ2\_M)} nel modello \textbf{Gemma-2-27b} mostra un leggero degrado delle prestazioni rispetto alla versione Q3\_K\_S, ma mantiene comunque correlazioni forti per la maggior parte dei tratti. Questo suggerisce un buon compromesso tra accuratezza e efficienza computazionale.

    \item Nei modelli \textbf{Gemma 2 9B} e \textbf{LLaMA 3.1 8B}, la differenza tra \textbf{Q8\_0} e \textbf{Q5\_K\_M} è minima, suggerendo che la quantizzazione a 5 bit potrebbe essere sufficiente per questi modelli più piccoli senza significative perdite di prestazioni. Questo risultato è particolarmente interessante per applicazioni che richiedono un equilibrio tra accuratezza e efficienza computazionale.
\end{itemize}
